% sec-04: Set B, applications 06 to 10.  Figure 19 (plantuml activity).
\section{Set B: The Medical Oncology Perspective}

\subsection{What Leads in Set B}

In the five medical oncology applications the drug and the patient selection
lead. Daraxonrasib (RMC-6236) is a RAS(ON) multi-selective inhibitor, given
perioperatively in KRAS G12-mutated resectable or borderline-resectable disease,
and the operation is described in full but as the instrumented setting that
makes the pharmacology measurable.

The endpoint table reorders accordingly: dose-limiting toxicity and the 3+3
escalation first, then pharmacokinetic and pharmacodynamic sampling timed around
the resection, then pathologic response, ctDNA clearance, and 24-month overall
survival. The operation matters to the pharmacology for one specific reason,
which no metastatic study can supply: it is the only setting in which paired
pre-treatment and on-treatment human tumour tissue can be obtained under
controlled timing.

\begin{apptable}
\begin{tabularx}{\textwidth}{>{\raggedright\arraybackslash}p{3.3cm} >{\raggedright\arraybackslash}p{3.3cm} >{\raggedright\arraybackslash}X}
\headrow \thead{Measurement} & \thead{Timing relative to resection} & \thead{Why no other setting supplies it}\\
\midrule
Plasma pharmacokinetics & Days $-14$, $-1$, and $+1$ & Exposure in a surgical population has not been characterised by any sponsor \\
Tumour pathway pharmacodynamics & Resection specimen, day 0 & Cannot be obtained in the metastatic setting at all \\
ctDNA clearance kinetics & Days $-14$, $+30$, $+90$ & Defines a shared surrogate anchored to a known resection date \\
Pathologic response grade & Resection specimen, day 0 & Anchors every later perioperative RAS study to one scale \\
\end{tabularx}
\end{apptable}
\tabcap{The four measurements the consortium in application 06 exists to make,
when each is taken relative to the resection, and the reason a metastatic trial
cannot supply any of them.}

\subsection{Applications 06 to 10}

\textbf{06, Foundation for the NIH.} Proposes a narrow pre-competitive question
that no single party owns: what happens pharmacologically when a RAS(ON)
inhibitor is given around a curative-intent pancreatic resection. Three parties
contribute, an independent scientist, an academic cancer center, and a drug
developer, and the single shared output, a specimen-level pharmacodynamic
dataset, is released openly and owned by none of them. That is what makes it
pre-competitive: it differentiates no product.

\textbf{07, HHMI Investigator Program.} Tests the one place the report's
person-based argument might not reach, an investigator with no academic
appointment, and asks the eligibility question directly rather than burying it.
The programme is written as a state machine rather than a set of aims, and its
defining property is that every failure transition re-enters the programme: if
exposure is inadequate the schedule changes and the agent is retained; if
pathway suppression is not demonstrated the agent changes and the question is
retained.

\textbf{08, NCI Cancer Therapy Evaluation Program.} The most clinically
technical of the ten, because a concept review is a technical review. It puts
one design point forward: a combined drug and device study in a surgical
population carries a hazard a standard escalation does not, in that an adverse
event may be attributable to the drug, to the operation, or to the advisory
system, and the dose decision depends on which. Attribution is therefore
adjudicated by a blinded reviewer before the toxicity window closes, and only
the drug branch reaches the dose decision.

\textbf{09, Convergent Research.} The only application whose organization states
its own end date. Three properties are put forward to be tested: the
organization closes in year five and the budget contains no year six; three
releases are scheduled of which only the specimen dataset depends on the trial
succeeding; and every artifact has one named recipient class fixed in year one,
because negotiating recipients in year five is how a dissolved organization
strands its own outputs.

\textbf{10, UC San Diego Moores Cancer Center.} Different in kind. It is
addressed to the partner institution rather than to a funder, asks for a
45-minute feasibility meeting rather than money, and carries three positioning
corrections drawn from the author's own site research: daraxonrasib is not
first-in-human; the robotic configuration must be named by manufacturer, model,
cleared configuration, arm count, instruments, and software version before
feasibility can be assessed; and no drug supply or letter of authorization
exists.

\subsection{What Happens If Any One Lands}

\begin{appfloat}[!tb]
\begin{appfig}[plantuml-type / activity with fork and join]
% Fork and join bars are 9.2cm wide, wider than the outermost lane, so they
% visibly span all four.  Four lanes at x = 0.4, 3.4, 6.4, 9.4; two boxes per
% lane at y = -1.0 and y = -2.3, giving 0.3 of clear space between lanes.
\node[umlinit] (st) at (4.9,1.5) {};
\node[umlkey,text width=44mm] (fund) at (4.9,0.75) {Any one of the ten applications is funded};
\draw[umlarrow] (st) -- (fund);
\node[umlbar,minimum width=9.2cm] (fk) at (4.9,-0.05) {};
\draw[umlarrow] (fund) -- (fk);
\foreach \i/\x/\a/\b in {1/0.4/{Execute the site agreement}/{Open the IRB submission},
  2/3.4/{Build the interlock rig}/{Measure stop latency on the bench},
  3/6.4/{Fix the advisory logging schema}/{Stand up the audit replay tool},
  4/9.4/{Request the drug cross-reference}/{Confirm pharmacy and supply chain}}{
  \node[umlbox,text width=24mm] (p\i) at (\x,-1.05) {\a};
  \node[umlbox,text width=24mm] (q\i) at (\x,-2.35) {\b};
  \draw[umlarrow] (\x,-0.11) -- (p\i.north);
  \draw[umlarrow] (p\i) -- (q\i);
  \draw[umlarrow] (q\i.south) -- (\x,-3.34);}
\node[umlbar,minimum width=9.2cm] (jn) at (4.9,-3.4) {};
\node[mmdec,text width=26mm] (d) at (4.9,-4.6) {Site agreement executed?};
\draw[umlarrow] (jn) -- (d);
\node[umlstateon,text width=32mm] (yes) at (1.6,-6.1) {Activate the IND, screen the first participant};
\node[umlstategray,text width=32mm] (no) at (8.2,-6.1) {Hold the cohort, re-scope to a second qualified site};
\draw[umlarrow] (d.south west) -- node[umlguard,pos=0.5] {[yes]} (yes.north east);
\draw[umlarrow] (d.south east) -- node[umlguard,pos=0.5] {[no]} (no.north west);
\node[umlfinal] (f1) at (1.6,-7.2) {};
\node[umlfinal] (f2) at (8.2,-7.2) {};
\draw[umlarrow] (yes) -- (f1);
\draw[umlarrow] (no) -- (f2);
\end{appfig}
\vspace{-0.7cm}
\figcaption{One award, four parallel workstreams, one join. The join condition
is the site agreement, not the money, which is why nine of the ten asks can be
satisfied by any single award rather than by all ten.}
\end{appfloat}

The join condition is worth stating plainly. Four workstreams run in parallel
once any one award lands, and three of them, the interlock rig, the logging
schema, and the pharmacy confirmation, complete without reference to which
funder paid. Only the site agreement gates what follows, and no funder can
supply it.

\subsection{The Two Sets Describe One Operation}

Set A and Set B are not two programmes. They describe the same hybrid procedure,
which carries surgical and medical oncology arms together: a staged eight-arm
robotic pancreaticoduodenectomy with perioperative daraxonrasib and an
on-premises advisory model, in up to 18 participants under a 3+3 escalation.

What differs is what each set leads with, and that difference is a claim about
reviewers rather than about science. A surgical reviewer asked to evaluate a
drug schedule will defer to the medical oncology literature; a medical oncology
reviewer asked to evaluate an anastomosis will defer to the surgical literature.
Writing the same trial twice, each time leading with what the reader can
actually assess, is the only honest way to ask both.
